% Generated draft: requires mathematical and source review.
\documentclass[11pt]{article}
\usepackage[margin=25mm]{geometry}
\usepackage{fontspec}
\setmainfont{Latin Modern Roman}
\usepackage{amsmath,amssymb,mathtools,unicode-math}
\setmathfont{Latin Modern Math}
\usepackage{enumitem,needspace,xurl,hyperref}
\setlength{\parindent}{0pt}
\setlength{\parskip}{5pt}
\setlength{\emergencystretch}{3em}
\newcommand{\N}{\mathbb{N}}
\newcommand{\Z}{\mathbb{Z}}
\newcommand{\Q}{\mathbb{Q}}
\newcommand{\R}{\mathbb{R}}
\newcommand{\C}{\mathbb{C}}
\newcommand{\F}{\mathbb{F}}
\newcommand{\A}{\mathbb{A}}
\newcommand{\PP}{\mathbb{P}}
\newcommand{\E}{\mathbb{E}}
\newcommand{\eps}{\varepsilon}
\newcommand{\dd}{\,\mathrm{d}}
\begin{document}
\section*{UnsolvedMath ID: 30005085}
\textbf{Source problem code:} catalogue entry
\subsection{Definitions and mathematical statement}
Let \(k\) be a field of characteristic \(p>0\), and let \(G\) be a reductive group with Weyl group \(W\). Write \(G_k^\vee\) for the dual reductive group over \(k\) intended by the source. Let \(X\) be its weight lattice and \(X^+\) the dominant weights relative to a choice of positive roots. The supplied excerpt does not specify the underlying root datum, the geometric base for \(G\), or the splitting assumptions needed to fix this setting completely. No further restriction on \(p\) is stated.

For \(\lambda\in X^+\), let \(L(\lambda)\) denote the irreducible rational \(G_k^\vee\)-module of highest weight \(\lambda\). For a subset \(\Omega\subseteq X^+\), define
\[
\mathcal R_\Omega
 =\bigl\langle L(\lambda)\mid\lambda\in\Omega\bigr\rangle_{\mathrm{Serre}}
 \subseteq\operatorname{Rep}(G_k^\vee),
\]
where the Serre subcategory generated by these modules is the smallest full subcategory closed under subobjects, quotients, and extensions and containing them. Here \(\operatorname{Rep}(G_k^\vee)\) is understood as the category of finite-dimensional rational representations; the supplied excerpt does not expressly specify this convention.

The source introduces the extended affine Weyl group
\[
W^{\mathrm{ext}}=W\ltimes X
\]
and describes \(\mathbin{\cdot}\) as its \(p\)-dilated dot action. However, both its displayed indexing condition and the clean statement use \(W\mathbin{\cdot}0\cap X^+\), with \(W\) previously defined as the finite Weyl group. Thus the two possible indexing sets are
\[
\Omega_{\mathrm{fin}}=(W\mathbin{\cdot}0)\cap X^+,
\qquad
\Omega_{\mathrm{ext}}=(W^{\mathrm{ext}}\mathbin{\cdot}0)\cap X^+.
\]
The source names the resulting category \(\operatorname{Rep}^{\mathrm{ext}}_0(G_k^\vee)\), but the supplied evidence does not resolve which orbit is intended. No identification of these two indexing sets is assumed here.

Let \(\operatorname{Gr}_G\) denote the affine Grassmannian of \(G\) in the geometric setting intended by the source, and let \(\mathrm{Iw}\) be an Iwahori subgroup of its loop group. Write
\[
\mathcal P=\operatorname{Perv}_{\mathrm{Iw}}(\operatorname{Gr}_G,k)
\]
for the category of perverse sheaves with coefficients in \(k\) that are constructible with respect to the stratification by \(\mathrm{Iw}\)-orbits. This notation follows the source's constructibility definition; it does not add an equivariance requirement. The geometric base, sheaf theory, and support conventions are not supplied.

\textbf{Conjectural target attributed to Finkelberg--Mirkovi\'{c}.}
After resolving the orbit convention and the omitted geometric hypotheses, does there exist an equivalence of categories
\[
Q:\mathcal P\xrightarrow{\sim}
\operatorname{Rep}^{\mathrm{ext}}_0(G_k^\vee)
\]
that geometrically realizes Frobenius twist? The source requires a commutative diagram, but that diagram is absent from the complete supplied record. Consequently, the required functors and compatibility isomorphism cannot be specified from this evidence, and the full conjecture is not yet a precise categorical statement.

The supplied account uses a Frobenius-twist functor \(\operatorname{Fr}\) with the asserted highest-weight rule
\[
\operatorname{Fr}(L(\lambda))\cong L(p\lambda)
\]
and asserts that these twists belong to \(\operatorname{Rep}^{\mathrm{ext}}_0(G_k^\vee)\). These assertions are retained as source claims. The original text says ``for every \(\lambda\in X\)'', although its highest-weight notation requires dominant weights; the intended quantifier and the Frobenius convention over the arbitrary field \(k\) require clarification.

For context, the excerpt describes geometric Satake via Tannakian reconstruction, identifying a reconstructed group with \(G_k^\vee\). It also states, under the additional hypothesis that \(k\) has good characteristic, that parity sheaves on \(\operatorname{Gr}_G\) are perverse and correspond under geometric Satake to tilting modules. These are supplied contextual claims, not additional hypotheses silently imposed on the conjecture or independently verified results.
\subsection{Short English statement}
The conjecture attributed to Finkelberg--Mirkovi\'{c} seeks an equivalence between Iwahori-orbit-constructible perverse sheaves on an affine Grassmannian and a specified Serre subcategory of representations of the dual group, compatible with Frobenius twist. The supplied record leaves the Weyl-group orbit ambiguous and omits the compatibility diagram, so the full statement requires clarification and human mathematical review.
\subsection{Sources}
\begin{itemize}
\item Local UnsolvedMath registry: \texttt{data/unsolvedmath/problems.json}, permanent ID \(30005085\), source problem code \texttt{catalogue entry}, title ``Finkelberg--Mirkovi\'{c} Equivalence in Modular Geometric Satake''. The supplied record attributes the conjecture to Finkelberg--Mirkovi\'{c}; it supplies no separate proposer.
\item Supplied report citation: \emph{Arbeitsgemeinschaft: Geometric Representation Theory} (2023), Oberwolfach Reports, pp.~929--978, DOI \texttt{10.4171/owr/2022/18}. \url{https://doi.org/10.4171/owr/2022/18}. The extracted passage carries page number 957. The record separately gives proposed year 2022; the relationship between these year labels is not resolved here.
\item Supplied literature source: Riche--Williamson, ``Smith--Treumann theory''. \url{https://doi.org/10.1007/s10240-022-00134-y}. The record supplies this abbreviated attribution and link, without a complete bibliographic entry.
\item Supplied literature source: ``Oberwolfach Report 2022/18''. \url{https://doi.org/10.4171/owr/2022/18}.
\item Literature assessment dated exactly \texttt{2026-08-22}: the supplied curation labels the problem open and states that the Finkelberg--Mirkovi\'{c} equivalence remains open in general, with substantial partial progress from geometric Satake and modular affine-Hecke advances. This is supplied dated curation, not an independently verified literature assessment.
\item The registry labels the statement \texttt{corrected\_verified} and reports a check against the official Oberwolfach report. This is supplied verification metadata; no source consultation or independent verification is claimed here. The excerpt's references [6] and [7] have no corresponding bibliographic entries in the supplied record.
\end{itemize}
\end{document}
